\section{Métodos clásicos de exploración}

\subsection{El planteamiento Multi-Armed Bandit}

\begin{knowledgebox}{Multi-Armed Bandit}
El problema bandit es la versión más simplificada del problema de exploración: hay $K$ brazos (acciones) y cada brazo tiene una distribución de recompensa desconocida $p(r | a)$. No hay transiciones de estado; el objetivo es maximizar la recompensa acumulada a lo largo de $T$ rondas de interacción. El arrepentimiento (regret) se define como la diferencia respecto al brazo óptimo:
$$\text{Regret}(T) = T \cdot \mu^* - \sum_{t=1}^T r_t$$
\end{knowledgebox}

\subsection{Cómo se define una «buena exploración» en un bandit}

El escenario bandit es importante porque en él todavía puede definirse con bastante claridad la «exploración óptima». El curso adopta el regret como métrica central: compara la recompensa acumulada que realmente obtienes con el rendimiento esperado de tirar siempre del brazo óptimo.

\begin{figure}[H]
\centering
\includegraphics[width=0.9\textwidth]{slides-images/slide-09.jpg}
\caption{Medida de optimalidad en un bandit: el regret, no la recompensa de una sola ronda}
\end{figure}

\subsection{$\epsilon$-Greedy}

La estrategia de exploración más sencilla: con probabilidad $1 - \epsilon$ se elige la acción que actualmente es la mejor y con probabilidad $\epsilon$ se elige una al azar.

\begin{warningbox}{Limitaciones de $\epsilon$-Greedy}
La exploración de $\epsilon$-Greedy es \textbf{no dirigida}: explora de manera aleatoria y uniforme sobre todas las acciones, sin considerar cuáles aportan más información. Cuando el espacio de acciones es grande o el entorno es complejo, esta exploración a ciegas resulta sumamente ineficiente.
\end{warningbox}

\subsection{Upper Confidence Bound (UCB)}

\begin{importantbox}{Algoritmo UCB}
La idea central de UCB es el \textbf{optimismo} (optimism in the face of uncertainty): asigna una estimación más alta a las acciones con mayor incertidumbre y, de ese modo, fomenta la exploración.

$$a_t = \arg\max_a \left[\hat{\mu}_a + c \sqrt{\frac{\ln t}{N_a}}\right]$$

donde $\hat{\mu}_a$ es la media empírica de la acción $a$ y $N_a$ es el número de veces que se ha elegido la acción $a$. El segundo término es una bonificación por incertidumbre: cuantas menos veces se ha elegido una acción, mayor es su incertidumbre y con mayor prioridad se explora.
\end{importantbox}

\subsection{Thompson Sampling}

\begin{importantbox}{Thompson Sampling}
Thompson Sampling adopta un enfoque bayesiano:
\begin{enumerate}
    \item Para la distribución de recompensa de cada acción se mantiene una distribución posterior $p(\mu_a | \text{data})$
    \item En cada ronda se muestrea de la posterior: $\tilde{\mu}_a \sim p(\mu_a | \text{data})$
    \item Se elige la acción con el valor muestreado más alto: $a_t = \arg\max_a \tilde{\mu}_a$
    \item Tras obtener la recompensa, se actualiza la posterior
\end{enumerate}
En la práctica, Thompson Sampling suele tener un desempeño excelente y, además, cuenta con un regret bound teóricamente óptimo.
\end{importantbox}

\begin{figure}[H]
\centering
\includegraphics[width=0.9\textwidth]{slides-images/slide-10.jpg}
\caption{Las dos líneas principales de la exploración en bandits: optimism under uncertainty y probability matching}
\end{figure}

\subsection{Por qué la teoría de bandits sigue siendo el punto de partida de la investigación sobre exploración}

El bandit parece demasiado simple, pero sigue siendo la vía de entrada clave para entender el problema de la exploración, porque es uno de los pocos planteamientos en los que todavía podemos ofrecer garantías de regret y discutir la «exploración óptima». Al pasar a MDP de horizonte largo, parcialmente observables y con estados de alta dimensión, muchas de estas conclusiones dejan de cumplirse directamente; por eso UCB y Thompson Sampling funcionan más como principios de diseño que se conservan que como plantillas de algoritmos que puedan copiarse tal cual.

\begin{knowledgebox}{Tres principios de diseño que se aprenden del bandit}
\begin{enumerate}
    \item La incertidumbre debe formar parte de la decisión; no basta con mirar la media empírica.
    \item La calidad de la exploración se juzga por el costo acumulado a largo plazo, no por el desempeño en un solo paso.
    \item En cuanto la tarea pasa de un 1-step bandit a una toma de decisiones de horizonte largo, la dificultad de la exploración aumenta drásticamente.
\end{enumerate}
\end{knowledgebox}

\subsection{Resumen de la sección}

Los métodos clásicos de exploración van desde el sencillo $\epsilon$-greedy hasta los métodos con fundamento de UCB y Thompson Sampling; la diferencia central entre ellos está en cómo aprovechan la incertidumbre para guiar la exploración.
